% ---
% id: reductions/cor:instances
% kind: cor
% title: "Three special cases of instrument-conditioned squashing"
% summary: >-
%   The branching rule applies to a convex mixture with a known weight, to a
%   measurement block diagonal for an orthogonal decomposition and to a
%   register read by a POVM. Only the known mixture fixes the branch weight as
%   a number without further information.
% prominence: standard
% status: accepted
% uses: [lem:master, lem:top]
% assumes: []
% refs: []
% provenance: newly-derived
% lean:
% ---

Three choices of the generalized instrument recur in this module's
constructions: a known mixture, a split of the input space into orthogonal
blocks, and a POVM measurement of a register that arrives with the light.

\begin{corollary}\label{cor:instances}
Let $\Lambda_j$ and $F^{(j)}$ be branch channels and branch targets as in
Lemma~\ref{lem:master}.
\begin{enumerate}[label=(\roman*),leftmargin=3.0em]
\item \emph{(convex mixture with a known weight)} If
  $M_v=(1-q)G_v+qR_v$ for POVMs $G,R$ on $\Bsys$ and $q\in[0,1]$, with
  $G\dsq_{\Lambda_G}F^{(G)}$ and $R$ \emph{arbitrary}, then
  $M\dsq\big\{F^{(G)}_v\dsum\dyad{v}_{\Flag}\big\}_v$, with the sound branch
  constraints $\tr[\Pi_1\sigma]=1-q$ and $\tr[\Pflag\sigma]=q$, in the operator
  form of Lemma~\ref{lem:master}\,(iv).
\item \emph{(block-wise squashing)} Let $\Hs{\Bsys}=\bigoplus_j\Hs{\Bsys_j}$ be
  a finite orthogonal decomposition and let every $M_v$ be block diagonal,
  $M_v=\bigoplus_jM^{(j)}_v$, with $M^{(j)}\dsq_{\Lambda_j}F^{(j)}$. Then
  $M\dsq\bigoplus_jF^{(j)}$.
\item \emph{(POVM-conditioned squashing)} If $\Bsys=\Bsys_0D$ with a POVM
  $\{N_j\}_{j\in J}$ on $D$ and $M_v=\sum_jM_{v|j}\otimes N_j$ with
  $M_{|j}\dsq_{\Lambda_j}F_{|j}$ on $\Qsys_j$, then
  $\{M_v\}_v\dsq_{\Lambda}\big\{\bigoplus_jF_{v|j}\big\}_v$ with
  \begin{equation}\label{eq:povmcond}
    \Lambda(\rho)=\bigoplus_{j\in J}\Lambda_j\Big(
      \tr_{D}\big[(\idon{\Bsys_0}\otimes N_j^{1/2})\,\rho\,
      (\idon{\Bsys_0}\otimes N_j^{1/2})\big]\Big).
  \end{equation}
\end{enumerate}
\end{corollary}

\begin{proof}
Apply Lemma~\ref{lem:master} with, respectively:

(i) $\Ins_1=(1-q)\idmap$ and $\Ins_2=q\,\idmap$, branch targets $F^{(G)}$ and
the flag PVM of Lemma~\ref{lem:top} (which exists for the unknown $R$), so that
$M_{v,1}=(1-q)G_v$ and $M_{v,2}=qR_v$ and part~(iii) gives
$\sum_jM_{v,j}=M_v$. Here $\Ins_j^{\dagger}(\idon{})$ is a multiple of
$\idon{\Bsys}$, which is the hypothesis of part~(iv).

(ii) $\Ins_j(\rho)\coloneqq P_j\rho P_j$ read as an operator on $\Bsys_j$, for
$P_j$ the projection onto $\Hs{\Bsys_j}$; then $\Ins_j^{\dagger}(X_j)$ is
$X_j$ embedded in the $j$-th block, so $M_{v,j}$ is $M^{(j)}_v$ so embedded and
part~(iii) gives $\sum_jM_{v,j}=M_v$. The instrument condition is
$\sum_jP_j=\idon{\Bsys}$.

(iii) $\Ins_j(\rho)\coloneqq
\tr_{D}\big[(\idon{\Bsys_0}\otimes N_j^{1/2})\rho(\idon{\Bsys_0}\otimes
N_j^{1/2})\big]$, which is CP as a Kraus map followed by a partial trace, and
satisfies $\sum_j\tr\Ins_j(\rho)=\sum_j\tr[(\idon{}\otimes N_j)\rho]=\tr\rho$.
Its adjoint is
$\Ins_j^{\dagger}(X)=(\idon{}\otimes N_j^{1/2})(X\otimes\idon{D})
(\idon{}\otimes N_j^{1/2})=X\otimes N_j$, the two factors commuting because
they act on different tensor factors; hence
$M_{v,j}=M_{v|j}\otimes N_j$ and part~(iii) gives
$\sum_jM_{v,j}=\sum_jM_{v|j}\otimes N_j=M_v$.
\end{proof}
